% FST-DE_DarkEnergy_Skeleton_v1_de.tex
% Entwurf -- Dunkle Energie als residuale Vakuum-Freie-Energie
% Status: ENTWURF (v1)
% Datum: 2026-03-10

\documentclass[12pt,a4paper]{article}
\usepackage[utf8]{inputenc}
\usepackage[T1]{fontenc}
\usepackage{lmodern}
\input{glyphtounicode}
\pdfgentounicode=1
\usepackage{amsmath,amssymb,amsthm}
\usepackage{array}
\usepackage{tabularx}
\usepackage{booktabs}
\usepackage{xcolor}
\usepackage{graphicx}
\usepackage{geometry}
\geometry{a4paper, margin=2.5cm}
\usepackage[ngerman]{babel}
\usepackage{xurl}
\usepackage{microtype}
\def\HyperDestNameFilter#1{de.#1}
\usepackage[unicode]{hyperref}
\hypersetup{
  colorlinks=true,
  linkcolor=blue!70!black,
  citecolor=green!50!black,
  urlcolor=blue!70!black,
  pdftitle={Dunkle Energie als residuale Vakuum-Freie-Energie: Eine thermodynamische Schranke für die kosmologische Konstante},
  pdfauthor={Lukas Geiger},
  pdfsubject={Preprint zu dunkler Energie und kosmologischer Konstante},
  pdfkeywords={dunkle Energie, kosmologische Konstante, Vakuumenergie, Skalar-Tensor-Gravitation, f(R)}
}
\setlength{\emergencystretch}{2em}
\newcolumntype{Y}{>{\raggedright\arraybackslash}X}

\theoremstyle{definition}
\newtheorem{definition}{Definition}[section]
\theoremstyle{plain}
\newtheorem{theorem}[definition]{Theorem}
\newtheorem{lemma}[definition]{Lemma}
\newtheorem{proposition}[definition]{Satz}
\newtheorem{corollary}[definition]{Korollar}
\newtheorem{axiom}[definition]{Axiom}
\theoremstyle{remark}
\newtheorem{remark}[definition]{Bemerkung}

% --- Eigene Befehle ---
\newcommand{\Leff}{\Lambda_{\mathrm{eff}}}
\newcommand{\LUV}{\Lambda_{\mathrm{UV}}}
\newcommand{\LIR}{\Lambda_{\mathrm{IR}}}
\newcommand{\rhov}{\rho_{\mathrm{vac}}}
\newcommand{\dd}{\mathrm{d}}

\title{{\Large\textbf{Dunkle Energie als residuale Vakuum-Freie-Energie:\\
Eine thermodynamische Schranke für die kosmologische Konstante}}\\
\large\textit{Teil des FST-Programms~\cite{GeigerFST-Hub-2026}}}

\author{Lukas Geiger\\
\small Unabhängiger Forscher, Bernau, Deutschland\\
\small ORCID: \href{https://orcid.org/0009-0005-7296-1534}{0009-0005-7296-1534}}

\date{Juli 2026}

\begin{document}
\maketitle
\vspace{-2.8em}

% ============================================================
\begin{abstract}\small
Diese Revision prüft den mathematischen Kern des vorgeschlagenen
thermodynamischen Dunkle-Energie-Rahmens. Für die dimensionsgeschlossene
reduzierte Funktion
\(\Phi_{\rm red}=A+B\rho\ln(\rho/\rho_{\rm Pl})+C\rho^2\), \(B,C>0\),
sind Existenz, Eindeutigkeit, strikte Konvexität und ein expliziter
Lambert-\(W\)-Minimierer bewiesen. Der ursprüngliche Impulsraumausdruck ist
dimensionsinhomogen und bestimmt die Matching-Koeffizienten nicht. Daher
bleibt die Skalierung
\(\Lambda_{\rm eff}\sim H_0^2/[c^2\ln(M_{\rm Pl}/H_0)]\) ein externer
RG-Matching-Ansatz. Ihr korrigierter Zahlenwert ist
\(3{,}8\times10^{-55}\,{\rm m}^{-2}\), etwa um den Faktor \(290\) kleiner
als der Beobachtungswert.

Auch der Skalarsektor erhält einen neuen Status. Das Pöschl--Teller-Potential
\(V_0\,{\rm sech}^2(\phi/\phi_0)\) besitzt bei \(\phi=0\) ein instabiles
Maximum; die frühere tanh-Historie ist ein phänomenologischer Proxy und kein
hergeleiteter stabiler Attraktor. Quasistatische Störungen in metrischem
Standard-\(f(R)\) nähern sich einem Slip von \(1/2\) (oder dem Kehrwert \(2\),
je nach Konvention) und einer ART-artigen Lensing-Kombination
\(\Sigma\simeq1/F\), nicht \(5/4\) und \(9/8\). Schließlich liefert der
Hu--Sawicki-Environment-Ledger korrekte notwendige Krümmungsschwellen, aber
kein source-bound lokales Profil. Weder
\(|f_{R0}|=1{,}2\times10^{-5}\) noch \(10^{-6}\) ist für sich ein
Cassini-Beweis. Das Rahmenwerk ist damit ein bedingtes Forschungsprogramm,
keine geschlossene Lösung des kosmologischen-Konstanten- oder
Screening-Problems.
\vspace{1.5em}
\noindent\textbf{Schlüsselwörter:} Kosmologische Konstante, Dunkle Energie,
Vakuumenergie, Skalar-Tensor-Gravitation, gravitativer Slip
\end{abstract}

\clearpage
% ============================================================
\section{Einleitung}\label{sec:intro}

Diese Revision dokumentiert den mathematischen Stand nach einer zeilenweisen
Beweis-/Paper-Prüfung am 16. Juli 2026. Ihr Zweck ist enger als der des
früheren Entwurfs: Sie trennt Aussagen, die aus einem dimensionsgeschlossenen
eindimensionalen Modell folgen, von physikalischen Identifikationen, die
offenbleiben.

Das Audit ändert vier Hauptaussagen. Erstens war der ursprüngliche
Impulsraum-Integrand dimensionsinhomogen und definierte ohne explizite
Matching-Koeffizienten keine physikalische Freie-Energie-Dichte. Zweitens
wurde das vorgeschlagene Pöschl--Teller-Sättigungsgesetz nicht aus dem
kanonischen Klein--Gordon--Friedmann-System hergeleitet; sein Hilltop ist
instabil. Drittens ist die als Ableitung einer Stationaritätsbedingung
definierte thermodynamische Beta-Funktion identisch null. Viertens sagt die
metrische Standard-\(f(R)\)-Störungstheorie nicht die früher angegebenen
Werte \(\eta=5/4\) und \(\Sigma=9/8\) voraus. Diese Claims werden im
Folgenden zurückgezogen.

\section{Dimensionsgeschlossenes Variationsmodell}\label{sec:functional}

\subsection{Rohe Motivation und ihre Grenze}

Der frühere motivierende Ausdruck lautete
\begin{equation}\label{eq:Phi-raw}
 \Phi_{\rm raw}(\rho)=
 \int_{\LIR}^{\LUV}
 \left[
 \frac{\omega(k)}{2}
 +T_{\rm dS}\rho\ln\!\left(\frac{\rho}{\rho_{\rm Pl}}\right)
 +\frac{8\pi G\rho^2}{k^2}
 \right]\frac{4\pi k^2}{(2\pi)^3}\,\dd k .
\end{equation}
In natürlichen Einheiten besitzen die drei Terme in der eckigen Klammer die
Dimensionen \(E\), \(E^5\) und \(E^4\). Sie können so nicht addiert werden.
Eine Planck-Einheiten-Skalierung repariert das Problem nicht, sondern verbirgt
Potenzen von \(M_{\rm Pl}\) und \(H\), also gerade die Potenzen, die die
behauptete kosmologische Skalierung bestimmen. Gleichung~\eqref{eq:Phi-raw}
ist daher nur eine schematische Motivation, kein physikalisches Funktional und
keine Herleitung aus einer Modensumme.

\subsection{Reduziertes Modell}

\begin{definition}[Reduzierte Freie-Energie-Dichte]\label{def:functional}
Für \(\rho>0\) sei
\begin{equation}\label{eq:Phi}
 \Phi_{\rm red}(\rho)
 =A+B\,\rho\ln\!\left(\frac{\rho}{\rho_{\rm Pl}}\right)+C\rho^2 ,
 \qquad B>0,\quad C>0 .
\end{equation}
Dabei gilt
\([\rho]=[\Phi_{\rm red}]=E^4\), \([A]=E^4\), \([B]=1\) und
\([C]=E^{-4}\). Die Konstanten \(B\) und \(C\) sind renormierte
Matching-Koeffizienten. Ihre Werte folgen nicht aus
Gl.~\eqref{eq:Phi-raw}.
\end{definition}

Der Kernel
\begin{equation}\label{eq:Vgrav}
 V_{\rm grav}(\rho,k)=\frac{8\pi G\rho^2}{k^2}
\end{equation}
hat die Dimension einer Energiedichte und kann nach einer explizit
spezifizierten Integration und Matching-Vorschrift einen positiven
quadratischen Term motivieren. Dimensionsanalyse wählt ihn nur dann als
Monom niedrigster Ordnung aus, wenn eine Potenz von \(G\), quadratische
Abhängigkeit von \(\rho\) und inverse Potenzen von \(k\) bereits vorausgesetzt
werden. Sie beweist keine Eindeutigkeit unter allgemeinen Kerneln; auch eine
quasistatische \(4/3\)-Kraftverstärkung leitet allein kein off-shell
Freie-Energie-Funktional her.

\section{Bewiesenes Resultat und bedingte Skalierung}\label{sec:main}

\begin{theorem}[Eindeutiger Minimierer des reduzierten Modells]\label{thm:main}
Die Funktion \(\Phi_{\rm red}\) aus Definition~\ref{def:functional} besitzt
genau einen globalen Minimierer \(\rho_\ast>0\). Er erfüllt
\begin{equation}\label{eq:stationarity}
 B\left[1+\ln\!\left(\frac{\rho_\ast}{\rho_{\rm Pl}}\right)\right]
 +2C\rho_\ast=0
\end{equation}
und ist explizit gegeben durch
\begin{equation}\label{eq:rho-lambert}
 \rho_\ast=\frac{B}{2C}\,
 W\!\left(\frac{2C\rho_{\rm Pl}}{eB}\right),
\end{equation}
wobei \(W\) die Hauptverzweigung der Lambert-Funktion ist.
\end{theorem}

\begin{proof}
Die erste Ableitung geht für \(\rho\to0^+\) gegen \(-\infty\) und für
\(\rho\to\infty\) gegen \(+\infty\). Außerdem gilt
\begin{equation}
 \Phi_{\rm red}''(\rho)=\frac{B}{\rho}+2C>0 .
\end{equation}
Damit besitzt die Ableitung genau eine Nullstelle; strikte Konvexität macht
sie zum eindeutigen globalen Minimierer. Umstellen von
Gl.~\eqref{eq:stationarity} ergibt
\((2C\rho_\ast/B)\exp(2C\rho_\ast/B)=2C\rho_{\rm Pl}/(eB)\) und damit
Gl.~\eqref{eq:rho-lambert}.
\end{proof}

Für die zugeordnete Krümmung verwenden wir die Konvention
\begin{equation}\label{eq:Leff}
 \Leff=\frac{8\pi G}{c^4}\rho_\ast .
\end{equation}
Eine Abhängigkeit von \(H_0\) folgt erst, wenn \(B/C\) unabhängig bestimmt
ist. Der häufig verwendete Ausdruck
\begin{equation}\label{eq:bound}
 \Leff^{\rm ansatz}
 \sim\frac{H_0^2}{c^2\ln(M_{\rm Pl}/H_0)}
\end{equation}
ist daher ein bedingter RG-Matching-Ansatz und keine Folge von
Theorem~\ref{thm:main}.

\begin{corollary}[Arithmetik des bedingten Ansatzes]\label{cor:numerical}
Für \(H_0=67{,}4\,{\rm km\,s^{-1}\,Mpc^{-1}}\) und
\(L=\ln(M_{\rm Pl}/H_0)\simeq140\) gilt
\begin{equation}
 \frac{H_0^2}{c^2L}\simeq3{,}8\times10^{-55}\ {\rm m}^{-2}.
\end{equation}
Verglichen mit
\(\Lambda_{\rm obs}\simeq1{,}11\times10^{-52}\ {\rm m}^{-2}\)
aus Planck~2018~\cite{Planck2018} liegt der Ansatz um etwa
\(2{,}9\times10^2\) zu niedrig. Das ist ein ungelöster
Normierungs-/Matching-Faktor und keine Übereinstimmung innerhalb einer
Größenordnung.
\end{corollary}

\subsection{Status des Skalarsektors}\label{sec:scalar-tensor}

Wir betrachten die phänomenologische Lagrangedichte
\begin{equation}\label{eq:Lagrangian}
 \mathcal L=
 \frac{R+\gamma R^2}{16\pi G}
 +\frac12(\partial\phi)^2
 -V_{\rm PT}(\phi)+\mathcal L_m ,
 \qquad
 V_{\rm PT}(\phi)=V_0\,{\rm sech}^2(\phi/\phi_0).
\end{equation}
Für den quadratischen Krümmungssektor gilt \(m_s^2=1/(6\gamma)\). Das
Skalarpotential erfüllt
\begin{equation}\label{eq:VPT-curvature}
 V_{\rm PT}'(0)=0,\qquad
 V_{\rm PT}''(0)=-\frac{2V_0}{\phi_0^2}<0 .
\end{equation}
Damit ist \(\phi=0\) ein Maximum, kein Minimum. Lineare Störungen erfüllen
\[
 r^2+3Hr-\frac{2V_0}{\phi_0^2}=0 ,
\]
und die größere Wurzel ist für jedes \(V_0>0\) positiv. Hubble-Reibung kann
das Rollen verzögern, den Hilltop aber nicht asymptotisch stabilisieren. Der
frühere Satz über einen stabilen de-Sitter-Endzustand und der
Oszillations-Forecast werden zurückgezogen.

Für ein kanonisches Skalarfeld und drucklose Materie definieren wir
\[
 x=\frac{\dot\phi}{\sqrt6\,HM_{\rm Pl}},\qquad
 y=\frac{\sqrt{V}}{\sqrt3\,HM_{\rm Pl}},\qquad
 \lambda=-M_{\rm Pl}\frac{V'}{V},\qquad N=\ln a .
\]
Das korrekte autonome System lautet
\begin{align}
 x'&=-3x+\sqrt{\frac32}\lambda y^2
      +\frac32x(1+x^2-y^2),\label{eq:dx}\\
 y'&=-\sqrt{\frac32}\lambda xy
      +\frac32y(1+x^2-y^2).\label{eq:dy}
\end{align}
Mit \(\Omega_\phi=x^2+y^2\) und
\(w_\phi=(x^2-y^2)/(x^2+y^2)\) gilt
\begin{equation}\label{eq:dOmega-exact}
 \frac{\dd\Omega_\phi}{\dd N}
 =-3w_\phi\,\Omega_\phi(1-\Omega_\phi).
\end{equation}
Dies reduziert sich nicht auf die in früheren Fassungen gedruckte
tanh-Gleichung.

Eine beschränkte tanh-Kurve darf weiterhin als Datenparametrisierung verwendet
werden, sofern sie so bezeichnet wird. Eine dimensionskonsistente,
nichtnegative Wahl ist
\begin{align}
 X_{\rm DE}^{\rm proxy}(a)
 &=\frac{\Phi_0}{2}\left[1+
 \tanh\!\bigl(\kappa(a-a_t)\bigr)\right],\label{eq:tanh-solution}\\
 \frac{\dd X_{\rm DE}^{\rm proxy}}{\dd a}
 &=2\kappa X_{\rm DE}^{\rm proxy}
 \left(1-\frac{X_{\rm DE}^{\rm proxy}}{\Phi_0}\right).
 \label{eq:saturation-ODE}
\end{align}
Hier ist \(X_{\rm DE}=\rho_{\rm DE}/\rho_{\rm crit,0}\), nicht der momentane
Anteil \(\rho_{\rm DE}/\rho_{\rm crit}(a)\). Die Gleichungen
\eqref{eq:tanh-solution}--\eqref{eq:saturation-ODE} sind ein
phänomenologischer Proxy und werden nicht aus Gl.~\eqref{eq:Lagrangian}
hergeleitet.

\section{Kinematische Folgen eines Dichte-Proxys}\label{sec:observations}

Die bedingte Abschätzung aus Gl.~\eqref{eq:bound} ist
\(3{,}8\times10^{-55}\,{\rm m}^{-2}\), nicht
\(3{,}8\times10^{-53}\,{\rm m}^{-2}\). Sie liegt daher etwa um den Faktor
\(290\) unter \(\Lambda_{\rm obs}\), nicht um den Faktor drei.

Eine beschränkte Kurve darf nur dann als phänomenologische Dichtehistorie
verwendet werden, wenn ihre Normierung konsistent definiert ist. Sei
\(X_{\rm DE}(a)=\rho_{\rm DE}(a)/\rho_{\rm crit,0}\) und
\[
 E^2(a)=\frac{H^2}{H_0^2}=\Omega_ma^{-3}+X_{\rm DE}(a).
\]
Dann gilt
\begin{align}
 q(a)&=-1-\frac{a}{2E^2}\frac{\dd E^2}{\dd a}
 =\frac{\tfrac12\Omega_ma^{-3}-X_{\rm DE}
 -\tfrac a2X_{\rm DE}'}{E^2},\label{eq:q-correct}\\
 w_{\rm DE}(a)&=-1-\frac13\frac{\dd\ln X_{\rm DE}}{\dd\ln a}.
 \label{eq:weff}
\end{align}
Die frühere Verzögerungsformel verwendete den falschen Koeffizienten für
\(aX_{\rm DE}'\). Außerdem wechselte sie zwischen einem momentanen
Dichteanteil und einer auf die heutige kritische Dichte normierten Größe.
Damit werden der berichtete Beschleunigungs-Redshift, \(w_0\), die
CPL-Abbildung und das Proxy-\(\chi^2\) bis zu einer Neuberechnung aus einer
physikalisch hergeleiteten Historie zurückgezogen. Eine offizielle
DESI-Likelihood-Auswertung wird nicht beansprucht.

\begin{remark}[DESI DR2 Phantom-Crossing und Skalar-Tensor-Rekonstruktionen]
Neuere Analysen von DESI-DR2-Daten in Kombination mit Planck 2018 und Supernovae Ia legen eine effektive Zustandsgleichung der Dunklen Energie $w(z)$ nahe, die die Phantom-Grenze $w=-1$ überquert. Während minimal gekoppelte einzelne Skalarfelder (Quintessenz) $w=-1$ nicht ohne Geist-Freiheitsgrade überqueren können, erzeugen rekonstruierte Skalar-Tensor-Gravitationstheorien~\cite{Efstratiou2025} und wechselwirkende Dunkle-Sektor-Modelle mit feldabhängigen Dunkle-Materie-Massen~\cite{Antusch2026,Sahoo2026} natürlicherweise ein scheinbares Phantom-Crossing $w_{\rm eff}(z) < -1$. Der phänomenologische Proxy in Gl.~\eqref{eq:tanh-solution} dient als kinematisches Testfeld für eine solche effektive Skalar-Tensor-Evolution.
\end{remark}

\section{Lokale Konsistenz-Gates}

\subsection{Strahlungszeitliche Spur}\label{sec:BBN}

Für ein ideales perfektes Fluid gilt in der hier verwendeten
Vorzeichenkonvention
\[
 T=(1-3w)\rho .
\]
Damit ist \(T=0\) für exakt masselose Strahlung mit \(w=1/3\). Dies entfernt
die führende klassische Spurquelle, ist aber kein Beweis der BBN-Sicherheit:
Massive Spezies, Spuranomalien, Hintergrundentwicklung und Anfangsbedingungen
müssen durch eine quantitative Häufigkeitsrechnung propagiert werden.

\subsection{Konvexität und Screening}\label{sec:screening-gamma}

Eine feste strikt konvexe Funktion von \(\rho\) besitzt einen Minimierer. Diese
elementare Tatsache ist kein No-Go-Satz für Chameleon-Screening. In einem
Chameleon-Modell hängt das effektive Skalarpotential von der Umgebungsdichte
ab; damit ändert sich die minimierte Funktion. Der frühere
\(\Gamma\)-Konvergenz- und Energie-Dissipations-Block lieferte keinen gültigen
Liminf-/Recovery-Beweis und vermischte Gradienten verschiedener Variablen; er
wird zurückgezogen. Die zugehörige numerische Phasengrafik bleibt höchstens
eine Toy-Illustration und keine Cassini-Rechnung.

\subsection{Reichweite der Variationsaussage}\label{sec:resolvent-vacuum}

Der Variationssatz betrifft nur die reduzierte Skalarfunktion mit vorgegebenen
positiven Koeffizienten. Er löst weder das Problem der radiativen Stabilität
noch bestimmt er ein Vakuum-Subtraktionsschema, konstruiert eine
RG-Trajektorie oder leitet ein lokales Screening-Profil her. Das sind
unabhängige Gates. Insbesondere ändern weder eine Resolventen-Analogie noch
eine Metrik auf dem Konfigurationsraum die Nullstellen des Differentials eines
festen Skalarfunktionals.

\section{Renormierung, Beta-Antwort und Beobachtungsstatus}
\label{sec:renorm}

Das physikalische Skalierungsproblem liegt in den Matching-Koeffizienten
\(B\) und \(C\). Eine gültige Herleitung muss ein quellen- und
schemakontrolliertes Ledger liefern:
\[
 {\rm bare}\longrightarrow{\rm Minkowski\text{-}Subtraktion}
 \longrightarrow{\rm Counterterms}
 \longrightarrow{\rm Horizont\text{-}Rest}.
\]
Eine solche Herleitung wird hier nicht geliefert. Die RG-, Datenverarbeitungs-,
Ein-Schleifen-, BBN- und Asymptotic-Safety-Passagen des älteren Entwurfs waren
Analogien oder Beobachtungen führender Ordnung, keine Beweise. Insbesondere
ist die Idealstrahlungs-Identität \(T=(1-3w)\rho=0\) für \(w=1/3\) korrekt,
garantiert aber allein keine BBN-Sicherheit im Sub-Prozent-Bereich, sobald
massive Spezies, Spuranomalien und die Hintergrund-Skalardynamik einbezogen
werden.

Die frühere Definition
\[
 \beta_\Lambda^{\rm alt}
 =\frac{\dd}{\dd\ln a}
 \left.\frac{\partial\Phi}{\partial\rho}\right|_{\rho=\rho_\ast(a)}
\]
ist identisch null, weil die Klammer für jedes \(a\) null ist. Sie kann keinen
nichtverschwindenden Fluss liefern. Differenziert man lediglich den bedingten
Ansatz \(\rho_\ast^{\rm ansatz}\propto H^2M_{\rm Pl}^2/L\) mit
\(L=\ln(M_{\rm Pl}/H)\), so ergibt sich
\begin{equation}\label{eq:beta-response}
 \frac{\dd\rho_\ast^{\rm ansatz}}{\dd\ln a}
 =\rho_\ast^{\rm ansatz}
 \left(2+\frac1L\right)\frac{\dd\ln H}{\dd\ln a}.
\end{equation}
Dies ist die Antwort eines angenommenen Skalierungsgesetzes, keine unabhängig
extrahierte Beta-Funktion und kein Beleg für das Skalierungsgesetz.

In dieser Revision wird weder eine Wilsonsche Beta-Funktion noch
Fixpunkt-Monotonie, eine Ein-Schleifen-Grenze für \(\gamma\) oder eine
Kullback--Leibler-Datenverarbeitungs-Aussage zum Resultat erhoben. Jede dieser
Aussagen benötigt eine dimensionsvollständige Wirkung, normalisierte
Wahrscheinlichkeitsobjekte für Informationsungleichungen und ein explizites
Renormierungsschema.

Die bedingte Skala
\[
 \Leff^{\rm ansatz}\simeq3{,}8\times10^{-55}\ {\rm m}^{-2}
\]
liegt etwa um den Faktor \(290\) unter dem Beobachtungswert. Dieses Paper
beansprucht daher weder eine Lösung des kosmologischen-Konstanten- oder
Koinzidenzproblems noch eine DESI-kompatible Zustandsgleichung oder einen
stabilen spätzeitlichen Attraktor.

\subsection{\texorpdfstring{Störungen in metrischem \(f(R)\)}{Störungen in metrischem f(R)}}\label{sec:predictions-falsifiable}

Wir verwenden die Konvention
\[
 \dd s^2=-(1+2\Psi)\dd t^2+a^2(1-2\Phi)\dd\mathbf x^2
\]
und definieren
\(s(k)= (k^2/a^2)/(k^2/a^2+m_s^2)\).
Für \(F=\dd f/\dd R\simeq1\) hat die quasistatische lineare Antwort die
Standardform
\begin{align}
 k^2\Psi&=-4\pi Ga^2\bar\rho\delta
 \left(1+\frac{s}{3}\right),\label{eq:Poisson-Psi}\\
 k^2\Phi&=-4\pi Ga^2\bar\rho\delta
 \left(1-\frac{s}{3}\right).\label{eq:Poisson-Phi}
\end{align}
Damit gilt bei der üblichen Zählerkonvention
\begin{equation}\label{eq:grav-slip}
 \eta_{\rm std}=\frac{\Phi}{\Psi}
 =\frac{1-s/3}{1+s/3}
 \longrightarrow\frac12
 \quad (s\to1),
\end{equation}
während die Kehrwertkonvention gegen \(2\) geht. Die Weyl-/Lensing-Kombination
bleibt bis auf den üblichen \(1/F\)-Faktor ART-artig:
\begin{equation}\label{eq:Sigma}
 \Sigma\simeq\frac1F\simeq1 .
\end{equation}
Folglich sind \(\eta=5/4\) und \(\Sigma=9/8\) keine Vorhersagen der
metrischen Standard-\(f(R)\)-Gravitation; alle darauf beruhenden Forecasts
werden zurückgezogen. Die \(4/3\)-Verstärkung betrifft das Kraftpotential
\(\Psi\), nicht einen universellen Freie-Energie-Koeffizienten
\cite{HuSawicki2007,SotiriouFaraoni2010}.

Die strikte Konvexität einer festen, umgebungsunabhängigen Skalarfunktion
\(\Phi_{\rm red}(\rho)\) schließt ein zweites Minimum derselben Funktion aus.
Sie schließt Chameleon-Screening nicht aus, weil sich dort das effektive
Skalarpotential mit der Umgebungsdichte ändert. Ebenso kann allein eine
positiv-definite Metrik auf dem Konfigurationsraum den Minimierer einer festen
Skalarfunktion nicht verschieben: \({\rm grad}_g\Phi=0\) gilt genau dann,
wenn \(\dd\Phi=0\).

\subsection{Source-bound lokales Gate}\label{sec:mcmc-results}

Der Environment-Ledger vom 2. Juli 2026 verwendet
\begin{equation}\label{eq:environment-contract}
 |f_R({\rm env})|
 =|f_{R0}|\left(\frac{R_0}{R_{\rm env}}\right)^{n+1}
\end{equation}
zusammen mit
\[
 \epsilon_{\rm th}
 =\frac{|f_{R,{\rm env}}-f_{R,{\rm in}}|}{2\Phi_\odot},
 \qquad f_{R,{\rm in}}\simeq0,\quad
 \Phi_\odot=2{,}12\times10^{-6}.
\]
Die Cassini-Schwelle \(\epsilon_{\rm th}<2{,}3\times10^{-5}\) erfordert
\begin{equation}
 \frac{|f_R({\rm env})|}{|f_{R0}|}
 <\frac{2\Phi_\odot(2{,}3\times10^{-5})}{|f_{R0}|}.
\end{equation}
Für \(n=1\) folgt
\[
 \frac{R_{\rm env}}{R_0}\ge350{,}787
 \quad\text{für }|f_{R0}|=1{,}2\times10^{-5},
 \qquad
 \frac{R_{\rm env}}{R_0}\ge101{,}264
 \quad\text{für }|f_{R0}|=10^{-6}.
\]
Diese Schwellenrechnungen sind arithmetisch korrekt. Derzeit liefert kein
unabhängiges lokales Dichte-/Krümmungsmodell das benötigte
\(R_{\rm env}/R_0\); der Ledger erlaubt daher kein Claim-Upgrade. Weder der
explorative Best-Fit \(1{,}2\times10^{-5}\) noch die konservative Leitplanke
\(10^{-6}\) ist für sich ein Cassini-Beweis. Der erste Wert ist nur
Hintergrund-Proxy-kompatibel und nicht nachgewiesen Cassini-sicher.

Der frühere explorative MCMC-Scan verwendete vereinfachte diagonale
Summary-Terme statt der offiziellen Planck-/DESI-Likelihoods und kodierte
Cassini durch einen linearen harten Proxy. Seine numerischen Samples dürfen
nur als Hintergrund-Stresstest gelesen werden; sie begründen weder
Modellpräferenz noch offizielle Parametergrenzen oder lokales Screening.

\section{Beweisstatus und nächste Schritte}\label{sec:discussion}

\begin{center}
\begin{tabularx}{\linewidth}{@{}lY@{}}
\toprule
\textbf{Punkt} & \textbf{Status nach dem Audit vom 16. Juli 2026}\\
\midrule
Reduziertes Funktional & Bewiesen: Existenz, Eindeutigkeit, strikte Konvexität und Lambert-\(W\)-Minimierer für vorgegebene \(B,C>0\).\\
Rohe Modensumme & Als dimensionsinhomogenes Funktional ungültig; nur Motivation.\\
Physikalische \(H_0^2/\ln\)-Skalierung & Offen und bedingt auf explizites RG-Matching; die korrigierte Rechnung verfehlt den Beobachtungswert um etwa \(290\).\\
Pöschl--Teller-Dynamik & Hilltop-Transient; kein stabiler de-Sitter-Attraktor und kein hergeleitetes tanh-Gesetz.\\
Metrischer \(f(R)\)-Slip/Lensing & Standardgrenzen sind \(\eta_{\rm std}\to1/2\) (oder Kehrwert \(2\)) und \(\Sigma\simeq1/F\), nicht \(5/4\) und \(9/8\).\\
Lokales Screening & Offen: Schwellenledger korrekt, aber alle numerischen Passes sind source-blocked.\\
Explorativer MCMC & Nur Hintergrund-Proxy; keine offizielle Planck-/DESI-Likelihood und keine Cassini-sichere Klassifikation.\\
\bottomrule
\end{tabularx}
\end{center}

Die nächsten entscheidenden Rechnungen sind: (i) dimensionsvolle,
renormierte \(B(H)\) und \(C(H)\) bestimmen; (ii) das Hilltop-Potential
ersetzen oder eine stabile Skalarhistorie aus den vollständigen Gleichungen
ableiten; (iii) ein unabhängiges galaktisches/solares Hu--Sawicki-Profil
liefern, dessen Krümmungsverhältnis die Ledger-Schwelle erfüllt; und (iv) erst
danach eine offizielle kosmologische Likelihood auswerten. Vor dem Schließen
dieser Gates ist weder ein Release noch ein Claim-Upgrade gerechtfertigt.

\clearpage

\subsection*{KI-Offenlegung}
Bei der Erstellung dieses Manuskripts wurden KI-gestützte Werkzeuge
(Claude, Anthropic) in unterstützender Funktion eingesetzt,
insbesondere für Literaturrecherche, Textstrukturierung und
Formulierungshilfe. Der konzeptionelle Entwurf, die wissenschaftliche
Argumentation und die Verantwortung für den gesamten Inhalt liegen
ausschließlich beim Autor.

\begin{thebibliography}{99}

\bibitem{Weinberg1989}
S.~Weinberg,
\textit{The cosmological constant problem},
Rev.\ Mod.\ Phys.\ \textbf{61}, 1--23 (1989).
\href{https://doi.org/10.1103/RevModPhys.61.1}{doi:10.1103/RevModPhys.61.1}.

\bibitem{Zeldovich1968}
Ya.~B.~Zeldovich,
\textit{The cosmological constant and the theory of elementary particles},
Sov.\ Phys.\ Usp.\ \textbf{11}, 381--393 (1968).
\href{https://doi.org/10.1070/PU1968v011n03ABEH003927}{doi:10.1070/PU1968v011n03ABEH003927}.

\bibitem{PeeblesRatra2003}
P.~J.~E.~Peebles and B.~Ratra,
\textit{The cosmological constant and dark energy},
Rev.\ Mod.\ Phys.\ \textbf{75}, 559--606 (2003).
\href{https://doi.org/10.1103/RevModPhys.75.559}{doi:10.1103/RevModPhys.75.559}.

\bibitem{Planck2018}
Planck Collaboration (N.~Aghanim et al.),
\textit{Planck 2018 results. VI. Cosmological parameters},
Astron.\ Astrophys.\ \textbf{641}, A6 (2020).
\href{https://doi.org/10.1051/0004-6361/201833910}{doi:10.1051/0004-6361/201833910}.

\bibitem{Carroll2001}
S.~M.~Carroll,
\textit{The cosmological constant},
Living Rev.\ Relativ.\ \textbf{4}, 1 (2001).
\href{https://doi.org/10.12942/lrr-2001-1}{doi:10.12942/lrr-2001-1}.

\bibitem{Riess1998}
A.~G.~Riess et al.,
\textit{Observational evidence from supernovae for an accelerating
universe and a cosmological constant},
Astron.\ J.\ \textbf{116}, 1009--1038 (1998).
\href{https://doi.org/10.1086/300499}{doi:10.1086/300499}.

\bibitem{Perlmutter1999}
S.~Perlmutter et al.,
\textit{Measurements of $\Omega$ and $\Lambda$ from 42 high-redshift
supernovae},
Astrophys.\ J.\ \textbf{517}, 565--586 (1999).
\href{https://doi.org/10.1086/307221}{doi:10.1086/307221}.

\bibitem{MartyushevSeleznev2006}
L.~M.~Martyushev and V.~D.~Seleznev,
\textit{Maximum entropy production principle in physics, chemistry and biology},
Phys.\ Rep.\ \textbf{426}, 1--45 (2006).
\href{https://doi.org/10.1016/j.physrep.2005.12.001}{doi:10.1016/j.physrep.2005.12.001}.

\bibitem{Starobinsky1980}
A.~A.~Starobinsky,
\textit{A new type of isotropic cosmological models without singularity},
Phys.\ Lett.\ B \textbf{91}, 99--102 (1980).
\href{https://doi.org/10.1016/0370-2693(80)90670-X}{doi:10.1016/0370-2693(80)90670-X}.

\bibitem{DESI2025}
DESI Collaboration (M.~Abdul-Karim et al.),
\textit{DESI DR2 Results II: Measurements of Baryon Acoustic Oscillations and Cosmological Constraints},
Phys.\ Rev.\ D \textbf{112}, 083515 (2025).
\href{https://doi.org/10.1103/tr6y-kpc6}{doi:10.1103/tr6y-kpc6}.
arXiv:2503.14738 [astro-ph.CO].

\bibitem{Eichhorn2020}
A.~Eichhorn,
\textit{Asymptotically safe gravity},
arXiv:2003.00044 (2020).

\bibitem{EichhornHeld2018}
A.~Eichhorn and A.~Held,
\textit{Top mass from asymptotic safety},
Phys.\ Lett.\ B \textbf{777}, 217--221 (2018).
\href{https://doi.org/10.1016/j.physletb.2017.12.040}{doi:10.1016/j.physletb.2017.12.040}.
arXiv:1707.01107.

\bibitem{HuSawicki2007}
W.~Hu and I.~Sawicki,
\textit{Models of $f(R)$ cosmic acceleration that evade solar-system tests},
Phys.\ Rev.\ D \textbf{76}, 064004 (2007).
\href{https://doi.org/10.1103/PhysRevD.76.064004}{doi:10.1103/PhysRevD.76.064004}.

\bibitem{SotiriouFaraoni2010}
T.~P.~Sotiriou and V.~Faraoni,
\textit{$f(R)$ theories of gravity},
Rev.\ Mod.\ Phys.\ \textbf{82}, 451--497 (2010).
\href{https://doi.org/10.1103/RevModPhys.82.451}{doi:10.1103/RevModPhys.82.451}.

\bibitem{BirrellDavies1982}
N.~D.~Birrell and P.~C.~W.~Davies,
\textit{Quantum Fields in Curved Space},
Cambridge Monographs on Mathematical Physics, Cambridge University Press, 1982.
\href{https://doi.org/10.1017/CBO9780511622632}{doi:10.1017/CBO9780511622632}.

\bibitem{MONDEntropy2025}
{\"O}.~Sevin{\c c}, {\"O}.~{\"O}kc{\"u}, and E.~Aydiner,
\textit{From MOND entropy to extended uncertainty principles: A unified framework},
Phys.\ Lett.\ B \textbf{873}, 140169 (2026).
\href{https://doi.org/10.1016/j.physletb.2026.140169}{doi:10.1016/j.physletb.2026.140169}.
arXiv:2601.14353 [gr-qc].

\bibitem{FinslerJCAP2025}
C.~Pfeifer, N.~Voicu, A.~Friedl-Sz\'asz, and E.~Popovici-Popescu,
\textit{From kinetic gases to an exponentially expanding universe -- the Finsler--Friedmann equation},
JCAP \textbf{2025}(10), 050 (2025).
\href{https://doi.org/10.1088/1475-7516/2025/10/050}{doi:10.1088/1475-7516/2025/10/050}.

\bibitem{HuntMOND2026}
T.~A.~Hunt,
\textit{Dark Matter Phenomenology as Informational Persistence: Nonlocal Gravity, MOND, and the Survival of Curvature Under Coarse-Graining},
Zenodo preprint (2026).
\href{https://doi.org/10.5281/zenodo.18215660}{doi:10.5281/zenodo.18215660}.

\bibitem{GeigerRH2026c}
L.~Geiger, \emph{From Landscape to Proof: The Even-Dominance Route to the Riemann Hypothesis}, Zenodo preprint, 2026.
\href{https://doi.org/10.5281/zenodo.19035640}{doi:10.5281/zenodo.19035640}.

\bibitem{FST-Unified2026}
L.~Geiger, \emph{FST Spectrum Duality: The Renormalized Free Energy Principle as Physical Instantiation of Functional Stability Theory}, Zenodo preprint, 2026.
\href{https://doi.org/10.5281/zenodo.19036190}{doi:10.5281/zenodo.19036190}.

\bibitem{FST-NS2026}
L.~Geiger, \emph{A Thermodynamic Obstruction to Finite-Time Blow-Up in the 3D Navier--Stokes Equations},
Zenodo preprint, 2026.
\href{https://doi.org/10.5281/zenodo.19087449}{doi:10.5281/zenodo.19087449}.

\bibitem{Geiger2026-YM}
L.~Geiger, \emph{Volume-Independent Spectral Gap for Strong-Coupling Lattice Gauge Theory via Poincar\'e--Dobrushin Machinery},
Zenodo preprint, 2026.
\href{https://doi.org/10.5281/zenodo.19087433}{doi:10.5281/zenodo.19087433}.

\bibitem{Geiger2026-BSD}
L.~Geiger, \emph{The BSD Conjecture as a Positivity Normal-Form Theorem},
Zenodo preprint, 2026.
\href{https://doi.org/10.5281/zenodo.19087443}{doi:10.5281/zenodo.19087443}.

\bibitem{CopelandLiddleWands1998}
E.~J.~Copeland, A.~R.~Liddle, and D.~Wands,
\textit{Exponential potentials and cosmological scaling solutions},
Phys.\ Rev.\ D \textbf{57}, 4686--4690 (1998).
\href{https://doi.org/10.1103/PhysRevD.57.4686}{doi:10.1103/PhysRevD.57.4686}.

\bibitem{Li2004}
M.~Li,
\textit{A model of holographic dark energy},
Phys.\ Lett.\ B \textbf{603}, 1--5 (2004).
\href{https://doi.org/10.1016/j.physletb.2004.10.014}{doi:10.1016/j.physletb.2004.10.014}.

\bibitem{BoussoPolchinski2000}
R.~Bousso and J.~Polchinski,
\textit{Quantization of four-form fluxes and dynamical neutralization of the cosmological constant},
JHEP \textbf{0006}, 006 (2000).
\href{https://doi.org/10.1088/1126-6708/2000/06/006}{doi:10.1088/1126-6708/2000/06/006}.

\bibitem{Susskind2003}
L.~Susskind,
\textit{The anthropic landscape of string theory},
in: \textit{Universe or Multiverse?}, Cambridge University Press, 247--266 (2007).
\href{https://doi.org/10.1017/CBO9781107050990.018}{doi:10.1017/CBO9781107050990.018}.

\bibitem{KaloperPadilla2014}
N.~Kaloper and A.~Padilla,
\textit{Sequestering the standard model vacuum energy},
Phys.\ Rev.\ Lett.\ \textbf{112}, 091304 (2014).
\href{https://doi.org/10.1103/PhysRevLett.112.091304}{doi:10.1103/PhysRevLett.112.091304}.

\bibitem{LinYang2004}
F.~Lin and Y.~Yang,
\textit{Existence of energy minimizers as stable knotted solitons in the Faddeev model},
Commun.\ Math.\ Phys.\ \textbf{249}, 273--303 (2004).
\href{https://doi.org/10.1007/s00220-004-1110-y}{doi:10.1007/s00220-004-1110-y}.

\bibitem{Wetterich2024}
C.~Wetterich,
\textit{Dark energy evolution from quantum gravity},
arXiv:2407.03465 [hep-th] (2024).

\bibitem{VasakCCGG2020}
D.~Vasak, J.~Struckmeier, and J.~Kirsch,
\textit{Dark energy and inflation invoked in CCGG by locally contorted space-time},
Eur.\ Phys.\ J.\ Plus \textbf{135}, 404 (2020).
\href{https://doi.org/10.1140/epjp/s13360-020-00415-7}{doi:10.1140/epjp/s13360-020-00415-7}.

\bibitem{ShapiroSola2009}
I.~L.~Shapiro and J.~Sol\`{a},
\textit{The scaling evolution of the cosmological constant},
JHEP \textbf{2002}(02), 006 (2002).
\href{https://doi.org/10.1088/1126-6708/2002/02/006}{doi:10.1088/1126-6708/2002/02/006}.
arXiv:hep-th/0012227.

\bibitem{SolaPeracaula2022}
J.~Sol\`{a} Peracaula,
\textit{The cosmological constant problem and running vacuum in the expanding universe},
Phil.\ Trans.\ R.\ Soc.\ A \textbf{380}, 20210182 (2022).
\href{https://doi.org/10.1098/rsta.2021.0182}{doi:10.1098/rsta.2021.0182}.

\bibitem{CohenKaplanNelson1999}
A.~G.~Cohen, D.~B.~Kaplan, and A.~E.~Nelson,
\textit{Effective Field Theory, Black Holes, and the Cosmological Constant},
Phys.\ Rev.\ Lett.\ \textbf{82}, 4971 (1999).
\href{https://doi.org/10.1103/PhysRevLett.82.4971}{doi:10.1103/PhysRevLett.82.4971}.

\bibitem{GeigerFST-Hub-2026}
L.~Geiger,
\textit{Functional Stability Theory --- A Programmatic Hub: Universal Convexity-Uniqueness Across Scales (From Particles to Cosmology)},
Zenodo (2026).
\href{https://doi.org/10.5281/zenodo.20130499}{doi:10.5281/zenodo.20130499}.

\bibitem{Efstratiou2025}
D.~Efstratiou, E.~A.~Paraskevas, and L.~Perivolaropoulos,
\textit{Addressing the DESI DR2 Phantom-Crossing Anomaly and Enhanced $H_0$ Tension with Reconstructed Scalar-Tensor Gravity},
arXiv:2511.04610 [astro-ph.CO] (2025).

\bibitem{Antusch2026}
S.~Antusch, S.~F.~King, and X.~Wang,
\textit{Coupled Dark Energy and Dark Matter for DESI: An Effective Guide to the Phantom Divide},
arXiv:2604.08449 [hep-ph] (2026).

\bibitem{Sahoo2026}
P.~Thanankullaphong, P.~Sahoo, P.~H.~Puttasiddappa, and N.~Roy,
\textit{Quintom Dark Energy: Future Attractor and Phantom Crossing in Light of DESI DR2 Observation},
arXiv:2601.02284 [gr-qc] (2026).

\end{thebibliography}

\end{document}
